\section{September 16, 2026}

\subsection{Comparing independent and spanning collections}

\begin{proposition}\label{prop:independent-versus-spanning}
  Suppose that $u_1,\ldots,u_m$ is a linearly independent collection in
  $V$ and that $w_1,\ldots,w_n$ spans $V$. Then
  \[
    m\leq n.
  \]
  Thus, in a finite-dimensional vector space, every linearly independent
  collection has cardinality at most that of every finite spanning
  collection.
\end{proposition}

\begin{proof}
  We successively insert the vectors $u_1,u_2,\ldots$ into the spanning
  list $w_1,\ldots,w_n$, removing one of the $w$'s after each insertion.

  First adjoin $u_1$ to obtain
  \[
    u_1,w_1,\ldots,w_n.
  \]
  This collection is linearly dependent because $w_1,\ldots,w_n$ spans
  $V$. Since $u_1\neq0$, one of the $w$'s is a linear combination of the
  other vectors. Remove that vector. The resulting collection still
  spans $V$ and has $n$ vectors.

  Inductively, suppose that after $k-1$ steps we have a spanning
  collection of $n$ vectors consisting of
  $u_1,\ldots,u_{k-1}$ and $n-k+1$ of the original $w$'s. Adjoin $u_k$.
  The resulting collection is dependent. In a nontrivial dependence
  relation, some remaining $w$ must have a nonzero coefficient;
  otherwise the relation would contradict the linear independence of
  $u_1,\ldots,u_k$. We can therefore remove that $w$ without changing
  the span.

  If $m>n$, then after $n$ steps the collection
  $u_1,\ldots,u_n$ would span $V$. In particular, $u_{n+1}$ would lie in
  their span, contradicting the linear independence of
  $u_1,\ldots,u_m$. Hence $m\leq n$.
\end{proof}

\begin{example}
  No collection of four vectors in $\bb{R}^3$ is linearly independent,
  because the three standard basis vectors span $\bb{R}^3$. Similarly,
  no collection of three vectors can span $\bb{R}^4$, because the four
  standard basis vectors are linearly independent in $\bb{R}^4$.
\end{example}

\subsection{Subspaces of finite-dimensional spaces}

\begin{proposition}
  Every subspace of a finite-dimensional vector space is
  finite-dimensional.
\end{proposition}

\begin{proof}
  Let $U$ be a subspace of a finite-dimensional vector space $V$. If
  $U=\{0\}$, then the empty collection spans $U$, so $U$ is
  finite-dimensional.

  Otherwise choose a nonzero vector $u_1\in U$. If
  $U=\operatorname{span}\{u_1\}$, we are done. If not, choose
  \[
    u_2\in U\setminus\operatorname{span}\{u_1\}.
  \]
  Continue in this fashion: once $u_1,\ldots,u_{k-1}$ have been chosen,
  stop if they span $U$; otherwise choose
  \[
    u_k\in U\setminus\operatorname{span}\{u_1,\ldots,u_{k-1}\}.
  \]
  At every stage the chosen collection is linearly independent. Since
  $V$ has a finite spanning collection, Proposition~
  \ref{prop:independent-versus-spanning} bounds the number of vectors in
  any linearly independent collection in $V$. The process must therefore
  stop after finitely many steps, at which point the chosen vectors span
  $U$.
\end{proof}

\begin{corollary}
  Every finite-dimensional vector space has a basis.
\end{corollary}

\begin{proof}
  By definition, a finite-dimensional vector space has a finite spanning
  collection. Removing redundant vectors from that collection produces
  a basis.
\end{proof}

\subsection{Extending a linearly independent collection}

\begin{proposition}[Basis extension theorem]
  Every finite linearly independent collection in a finite-dimensional
  vector space can be extended to a basis of the space.
\end{proposition}

\begin{proof}
  Let $u_1,\ldots,u_m$ be linearly independent in $V$, and choose a
  finite spanning collection $w_1,\ldots,w_n$ for $V$. The combined
  collection
  \[
    u_1,\ldots,u_m,w_1,\ldots,w_n
  \]
  spans $V$. Examine the $w$'s in order. Whenever a vector $w_j$ lies in
  the span of the vectors preceding it, remove it. Removing such a
  vector does not change the span, and none of the $u_i$ is removed.
  After all the $w$'s have been considered, the remaining collection
  spans $V$ and is linearly independent. It is therefore a basis
  containing $u_1,\ldots,u_m$.
\end{proof}

\begin{example}
  Consider the linearly independent vectors
  \[
    v_1=\begin{bmatrix}-2\\3\\4\end{bmatrix},
    \qquad
    v_2=\begin{bmatrix}-9\\6\\8\end{bmatrix}
  \]
  in $\bb{R}^3$. Adjoin the standard basis vectors
  $e_1,e_2,e_3$. Since
  \[
    e_1=-\frac15(v_2-2v_1),
  \]
  the vector $e_1$ is redundant. On the other hand, $e_2$ is not in
  $\operatorname{span}\{v_1,v_2\}$, so we retain it. Finally,
  \[
    e_3=-\frac1{10}\left(v_2-\frac92v_1+\frac{15}{2}e_2\right),
  \]
  so $e_3$ is redundant. Therefore
  \[
    \left\{
      \begin{bmatrix}-2\\3\\4\end{bmatrix},
      \begin{bmatrix}-9\\6\\8\end{bmatrix},
      \begin{bmatrix}0\\1\\0\end{bmatrix}
    \right\}
  \]
  is a basis of $\bb{R}^3$ extending the original collection.
\end{example}
